\chapter{Authoring Guide}
\label{ch:authoring-guide}

This chapter teaches you to write a \texttt{.home.yaml} property file from scratch. It walks
through zone and asset declaration, routine authoring, incident templates, delegation
configuration, and consumable tracking, ending with a complete pool-maintenance example.

% ─────────────────────────────────────────────────────────────────────────────
\section{Starting a New Property File}
\label{sec:authoring-start}
% ─────────────────────────────────────────────────────────────────────────────

Every \texttt{gert.home} property file begins with the Kit declaration and a minimal
\texttt{property} block.

\begin{minted}{yaml}
apiVersion: runbook/v1
kit: home/v0

property:
  name: My Home
  zones: []

assets: []
routines: []
incident_templates: []
consumables: []
delegations: []
\end{minted}

Start by naming the property and identifying the major areas (zones) that require maintenance.
Assets, routines, and incidents can be added incrementally. A property file with only zones
and assets is valid and will compile without error.

% ─────────────────────────────────────────────────────────────────────────────
\section{Declaring Zones and Assets}
\label{sec:authoring-zones-assets}
% ─────────────────────────────────────────────────────────────────────────────

\paragraph{Zones} are named areas of the property. Use zones for locations that contain
multiple assets or are referenced by multiple routines (e.g., \texttt{pool},
\texttt{front\_lawn}, \texttt{garage}).

\paragraph{Assets} are specific tracked physical items. Use assets when you need to record
installation date, warranty expiry, or model metadata for a particular device or appliance.

\begin{minted}{yaml}
property:
  name: Casa Santiago
  zones:
    - id: pool
      label: Swimming Pool
      description: Saltwater pool with heat pump and LED lighting
    - id: front_lawn
      label: Front Lawn

assets:
  - id: pool_pump
    zone: pool
    label: Hayward SuperPump VS
    installed: 2023-06-15
    warranty_expires: 2025-06-15
    metadata:
      model: VS-100
\end{minted}

\begin{tcolorbox}[colback=yellow!10,colframe=orange!70,title=Rule of Thumb]
  Use a \textbf{zone} for any area you refer to by location (``the pool area'', ``the
  garage''). Use an \textbf{asset} for any specific item you might need to look up by model
  number, warranty, or service history.
\end{tcolorbox}

% ─────────────────────────────────────────────────────────────────────────────
\section{Defining Routines}
\label{sec:authoring-routines}
% ─────────────────────────────────────────────────────────────────────────────

Routines are the core of the property file. Each routine declares a recurring task with a
cadence, an evidence requirement, and optional notifications.

\begin{minted}{yaml}
routines:
  - id: pool_clean
    label: Pool cleaning
    zone: pool
    cadence:
      every: 7d
    evidence:
      type: photo
      prompt: Take photo of clear water and chemical test strip
    notifications:
      remind_hours_before: 4
      escalate_hours_overdue: 24

  - id: lawn_mow
    label: Mow front lawn
    zone: front_lawn
    cadence:
      seasonal:
        summer: 7d
        autumn: 10d
        winter: 21d
        spring: 10d
    evidence:
      type: photo
    executor:
      role: gardener
\end{minted}

\paragraph{Fixed vs.\ seasonal cadence.}
Use \texttt{every: Nd} when the task frequency does not change with the seasons (e.g., pool
cleaning, bin collection). Use \texttt{seasonal} when frequency varies significantly across
the year (e.g., lawn mowing, irrigation). Mixing both cadence types in the same routine
is not allowed --- choose one.

% ─────────────────────────────────────────────────────────────────────────────
\section{Using \texttt{home.incident}}
\label{sec:authoring-incidents}
% ─────────────────────────────────────────────────────────────────────────────

Use \texttt{incident\_templates} for reactive, event-driven workflows --- situations you
cannot schedule in advance. Common examples: water leaks, appliance failures, power outages.

\begin{minted}{yaml}
incident_templates:
  - id: leak
    label: Water leak
    zones: [pool, backyard, garage]
    steps:
      - id: assess
        type: human_task
        label: Locate and assess the leak
        evidence:
          type: photo
          prompt: Take photo of leak source and affected area
      - id: shutoff
        type: human_task
        label: Shut off water supply
        depends_on: [assess]
        evidence:
          type: photo
      - id: severity_check
        type: decision
        label: Can you fix it yourself?
        depends_on: [shutoff]
        choices:
          - value: yes
            label: Yes, I can fix it
            next_step: fix
          - value: no
            label: No, call a plumber
            next_step: call_plumber
      - id: fix
        type: human_task
        label: Perform repair
        evidence:
          type: photo
      - id: call_plumber
        type: human_task
        label: Contact a licensed plumber
        evidence:
          type: note
          prompt: Record plumber name, contact number, and estimated arrival time
\end{minted}

\paragraph{Structuring steps.}
Use \texttt{depends\_on} to chain steps sequentially. Steps without \texttt{depends\_on}
may start immediately when the incident is opened. Use \texttt{type: decision} for branching
choices and \texttt{type: parallel} for tasks that can be done concurrently.

% ─────────────────────────────────────────────────────────────────────────────
\section{Configuring Delegation}
\label{sec:authoring-delegation}
% ─────────────────────────────────────────────────────────────────────────────

Add a \texttt{delegations} entry before you leave the property. Specify the delegate's
contact details, the exact date range, which routines or zones they are assigned to, and
what permissions they have.

\begin{minted}{yaml}
delegations:
  - delegate:
      name: Carlos Méndez
      contact: "+56 9 8765 4321"
      email: carlos@example.com
    active:
      from: 2025-04-25
      to: 2025-05-02
    assigns:
      - routine: water_plants
      - zone: pool
    permissions:
      can_report_incidents: true
      can_modify_routines: false
      can_view_history: true
    notifications:
      remind_delegate_hours_before: 24
      notify_owner_on_completion: true
\end{minted}

\paragraph{Timing the away window.}
Set \texttt{active.from} to the day you leave and \texttt{active.to} to the day you return.
The delegation activates at midnight on \texttt{active.from} and expires at midnight on
\texttt{active.to} (in the timezone set by \texttt{GERT\_HOME\_TZ}).

\paragraph{Scoping assigns.}
Use \texttt{\{routine: id\}} to assign specific routines, and \texttt{\{zone: id\}} to
assign all routines whose \texttt{zone} field matches. Zone assignment is useful when you want
the delegate to handle everything in a given area without listing each routine individually.

% ─────────────────────────────────────────────────────────────────────────────
\section{Tracking Consumables}
\label{sec:authoring-consumables}
% ─────────────────────────────────────────────────────────────────────────────

Consumables track items that must be replaced on a schedule. The runtime computes the next
replacement date from \texttt{last\_replaced} + \texttt{replace\_every} and alerts when due.

\begin{minted}{yaml}
consumables:
  - id: pool_filter_sand
    label: Pool sand filter media
    zone: pool
    asset: pool_pump
    replace_every: 180d
    last_replaced: 2024-11-01
    stock:
      current: 1
      reorder_at: 1
    triggers_routine: pool_filter_replace
\end{minted}

After replacing a consumable, update \texttt{last\_replaced} to today's date. If
\texttt{triggers\_routine} is set, the runtime will fire the named routine when the
consumable's replacement date arrives.

% ─────────────────────────────────────────────────────────────────────────────
\section{Evidence Best Practices}
\label{sec:authoring-evidence}
% ─────────────────────────────────────────────────────────────────────────────

\begin{itemize}
  \item Use \texttt{photo} for outdoor and physical tasks where completion is visible
        (mowing, pool cleaning, leak repair). Photos provide tamper-evident proof and are
        SHA-256 hashed in the run record.
  \item Use \texttt{note} for tasks involving cost, measurements, or contractor
        communications (e.g., recording chemical readings, logging repair costs).
  \item Use \texttt{checklist} for multi-step verification tasks where partial completion
        must be tracked (e.g., pre-departure safety checks, appliance service checklists).
  \item Use \texttt{none} sparingly --- only for tasks where completion is self-evident and
        no record is required (e.g., resetting a timer, clearing a filter alarm).
  \item Always write a \texttt{prompt} for \texttt{photo} and \texttt{note} evidence. A
        clear prompt reduces ambiguous submissions and improves the audit trail.
  \item For incident steps, require \texttt{photo} evidence on the final \texttt{human\_task}
        step so you have documented proof of the completed repair.
\end{itemize}

% ─────────────────────────────────────────────────────────────────────────────
\section{Complete Example: Pool Maintenance Property}
\label{sec:authoring-complete-example}
% ─────────────────────────────────────────────────────────────────────────────

The following is a complete \texttt{.home.yaml} file for a property with a swimming pool and
front lawn. It demonstrates all major features: zones, assets, fixed and seasonal routines,
an incident template, a consumable, and a delegation.

\begin{minted}{yaml}
apiVersion: runbook/v1
kit: home/v0

property:
  name: Casa Santiago
  zones:
    - id: pool
      label: Swimming Pool
      description: Saltwater pool with heat pump and LED lighting
    - id: front_lawn
      label: Front Lawn

assets:
  - id: pool_pump
    zone: pool
    label: Hayward SuperPump VS
    installed: 2023-06-15
    warranty_expires: 2025-06-15
    metadata:
      model: VS-100

routines:
  - id: pool_clean
    label: Pool cleaning
    zone: pool
    cadence:
      every: 7d
    evidence:
      type: photo
      prompt: Take photo of clear water and chemical test strip
    notifications:
      remind_hours_before: 4
      escalate_hours_overdue: 24

  - id: lawn_mow
    label: Mow front lawn
    zone: front_lawn
    cadence:
      seasonal:
        summer: 7d
        autumn: 10d
        winter: 21d
        spring: 10d
    evidence:
      type: photo
    executor:
      role: gardener

incident_templates:
  - id: leak
    label: Water leak
    zones: [pool, backyard, garage]
    steps:
      - id: assess
        type: human_task
        label: Locate and assess the leak
        evidence:
          type: photo
          prompt: Take photo of leak source and affected area
      - id: shutoff
        type: human_task
        label: Shut off water supply
        depends_on: [assess]
        evidence:
          type: photo
      - id: fix
        type: human_task
        label: Perform repair
        depends_on: [shutoff]
        evidence:
          type: photo

consumables:
  - id: pool_filter_sand
    label: Pool sand filter media
    zone: pool
    asset: pool_pump
    replace_every: 180d
    last_replaced: 2024-11-01
    stock:
      current: 1
      reorder_at: 1

delegations:
  - delegate:
      name: Carlos Méndez
      contact: "+56 9 8765 4321"
    active:
      from: 2025-04-25
      to: 2025-05-02
    assigns:
      - routine: pool_clean
      - zone: front_lawn
    permissions:
      can_report_incidents: true
      can_modify_routines: false
      can_view_history: true
    notifications:
      remind_delegate_hours_before: 24
      notify_owner_on_completion: true
\end{minted}
